% One-time imported draft. Editable native source; no Markdown regeneration.
\section{Step 2 Connect the model and locate its
  port}\label{005-seven-segment-counter-s02-step-2-connect-the-model-and-locate-its-port}

The source-header comment configures the plugin tool. The .eqv
declaration gives a source-level constant a name; it reserves no memory.
\texttt{RIGHT\_DISPLAY} is the simulated device address used by the
program.

\subsection{Write or
  inspect}\label{005-seven-segment-counter-s02-write-or-inspect}

Inspect the supplied header, constant, table, and input words. Keep the
Default layout and little-endian setting for this lesson.

\subsection{Expected
  observation}\label{005-seven-segment-counter-s02-expected-observation}

Digital Lab Sim should be connected for this execution. Its right port
is 0xffff0010; its neighboring control ports are not part of this
counter.

\subsection{Check your
  understanding}\label{005-seven-segment-counter-s02-check-your-understanding}

Does .eqv \texttt{RIGHT\_DISPLAY} allocate a byte for the display?

\begin{itemize}
  \tightlist
  \item
        \textbf{A} Yes; it creates the peripheral in ordinary data memory.
  \item
        \textbf{B} Yes; it also writes the first glyph.
  \item
        \textbf{C} No; it names a constant used during assembly.
\end{itemize}
